\section{Choosing an active space}
\label{sec:active-space}

The smallest useful active space for a vertical excitation is the
frontier pair: the highest occupied orbital and the lowest unoccupied
one, two electrons in two orbitals (or more, when the window must
widen; see below).  This chapter is about the cases where that default
silently selects the wrong orbitals, how to detect them, and what to
do instead.  The detection rules live in
\texttt{diffmrc.analysis.orbitals} and cost one SCF to apply.

\subsection{The default window, and the cases that break it}

The default window takes the frontier pair in energy order.  It is
correct for the common chromophores: the carbonyl $n \to \pi^*$ pair
of aldehydes and ketones, and the $\pi \to \pi^*$ pair of a simple
alkene.  Widening it is unnecessary there, and this tutorial's
benchmark sets it as the baseline.

Two patterns break it.

\textbf{A linear molecule with two pi systems.}  In a nitrile, a
diyne, or cyanogen, the occupied pi orbitals and the virtual pi
orbitals each come as degenerate perpendicular components ($x$ and
$y$).  An energy-ordered window can pair the $x$ component of the
occupied set with the $y$ component of the virtual set.  That
excitation is symmetry-forbidden, its exchange integral vanishes, and
the computed excitation energy is systematically off --- measured at
up to 2~eV for the molecules in this tutorial's benchmark set.

\textbf{An amide.}  In formamide the conjugated $\pi/\pi^*$ pair
sits on top of the frontier, while the lone pair that gives the
characteristic $n \to \pi^*$ transition lies one orbital lower.  The
default window computes a $\pi \to \pi^*$ state and compares it to an
$n \to \pi^*$ reference: a mismatch of more than 2~eV in the state
that was supposed to be reproduced.

\subsection{Auditing a window in one SCF}

The audit is a per-axis decomposition of the frontier orbitals' $p$
weight.  A delocalized single-axis orbital is a pi; an atom-centred
single-axis orbital is a lone pair; a window whose HOMO and LUMO sit
on different axes while the occupied frontier is degenerate is a
cross-component pair.

\begin{lstlisting}[language=Python]
import numpy as np
from diffmrc.analysis.orbitals import heavy_atoms, frontier_report

# mol and mf from a plain pyscfad SCF (see the first-calculation chapter)
labels = mol.ao_labels()
C = np.asarray(mf.mo_coeff)
e = np.asarray(mf.mo_energy)
nocc = mol.nelectron // 2

rep = frontier_report(C, labels, e, nocc, heavy_atoms(labels),
                      target='npi', nelectron=mol.nelectron)
print(rep['homo']['kind'], rep['cross_component'],
      rep['suggested_window'])
\end{lstlisting}

\texttt{frontier\_\allowbreak report} returns everything that entered the
suggestion --- the HOMO-1/HOMO/LUMO identities, the degeneracy map of
the occupied frontier, the index of a lone pair below the HOMO when
the target needs one --- so that the suggestion can be checked rather
than trusted.  For formamide it reports that the HOMO is a
delocalized pi, finds the lone pair two orbitals below, and suggests
a three-orbital window spanning lone pair, pi, and pi$^*$.

\subsection{The corrected windows}

For the two failure modes the recipe is the same three-orbital
frontier window, three orbitals instead of two:

\begin{itemize}
  \item \textbf{Linear two-pi molecules}: start the window one orbital
    lower than the default, so that both components of each degenerate
    pair are inside and the CI can pick the in-component combinations.
  \item \textbf{Amides}: start the window at the lone pair, three
    orbitals wide, so that the lone pair, the pi system, and
    pi$^*$ are all active.
\end{itemize}

One bookkeeping rule travels with the wider window: the electron
count must follow the window, not be hard-coded.  The scripts in this
project use

\begin{equation}
  n_\mathrm{ele} = N_\mathrm{ele} - 2\, n_\mathrm{core},
\end{equation}

which is what the three-orbital windows above need (four active
electrons for the linear two-pi molecules, not two).

The effect on the numbers, from this project's benchmark records:

\begin{center}
\begin{tabular}{lccc}
\toprule
System & State & default window & corrected window \\
\midrule
Cyanogen, $\Sigma_u^-$ singlet & $\pi\pi^*$ & $+1.05$ & $+0.62$ \\
Cyanogen, $\Sigma_u^+$ triplet & $\pi\pi^*$ & $+1.95$ & $+1.20$ \\
Formamide, $A''$ singlet & $n\pi^*$ & $+2.12$ & $+0.11$ \\
\bottomrule
\end{tabular}
\end{center}

The deviations are against the reference set's best estimates (eV).
The corrected formamide window reproduces the reference
\textsc{sc}-NEVPT2 value to $0.02$~eV.  The residual deviations of the
linear two-pi systems are a separate story --- they survive every
window fix and are discussed in Chapter~\ref{sec:limitations}.

\subsection{When no element rule applies}
\label{sec:fallback-selection}

The rules above are written for the lanthanide families.  A main-group
molecule used to be refused outright, and a benchmark of six systems
found the release selecting only one of them.  The driver now falls
back to the occupation criterion, which is the same signal the
lanthanide path reads:

\begin{itemize}
\item every orbital whose reference occupation lies inside
  $(0.02,\,1.98)$ enters the space;
\item a closed-shell reference has no such orbital, so the frontier
  pair --- highest occupied and lowest virtual --- is taken instead;
\item a degenerate partner of any selected orbital enters with it,
  because half a degenerate pair describes a symmetry-broken state.
\end{itemize}

On the benchmark set this turns all four main-group cases into
selections that converge at the energy and gradient level:
H$_2$O takes orbitals 4 and 5, triplet CH$_2$ takes its two singly
occupied orbitals, N$_2$ takes the completed space $\{5,6,7,8\}$ (with
the curvature fallback described below), and O$_3$ takes the frontier
pair on its C$_{2v}$ equilibrium geometry, with the CASSCF energy
0.114~Eh below the RHF reference.

Converging is not the same as being a minimum, and the curvature check
is a separate gate.  Of these, the triplet CH$_2$ and N$_2$ spaces are
genuine minima (lowest orbital-Hessian eigenvalues $+0.555$ and
$+0.065$).  The H$_2$O frontier pair has three negative eigenvalues,
and the three-orbital space of the equilateral D$_{3h}$ ozone ring has
thirteen, because the occupied $e'$ pair $0.137$~Ha below the frontier
is frozen in the core and the orbitals relax along that direction.
Both are reported as refusals with the curvature evidence.  The ozone
selection is also geometry dependent: the C$_{2v}$ equilibrium geometry
gives a clean frontier pair, and the ring geometry is refused.

The retry is decided by curvature, not by the gradient alone.  A
completed space can converge at both the energy and the gradient level
and still be an orbital saddle.  The driver therefore evaluates the
curvature of the proposed space and, when the completion came from the
occupation fallback and the space holds more than two orbitals, of the
frontier pair as well, then keeps the candidate with fewer negative
eigenvalues.  The equilateral D$_{3h}$ ozone ring is the measured case:
the three-orbital completion has thirteen negative eigenvalues and the
frontier pair two, so the frontier pair is kept and the run is refused
with both spectra as evidence.  The H$_2$O frontier pair has three
negative eigenvalues and is refused in the same way.  N$_2$/6-31G's
completed space $\{5,6,7,8\}$ is a minimum with zero negative
eigenvalues, and the space itself is kept.  Its analytic orbital Hessian
is not finite at the symmetric stationary point, so the refinement
rebuilds the curvature from symmetric differences of the analytic
gradient; two steps take the gradient from $7.3\times10^{-5}$ to
$8.8\times10^{-14}$ and the space passes the robust level.  The retry
is automatic only for a space the fallback proposed; an active space
given by hand is used as given.

A manual list bypasses the degenerate-partner rule, and a window that
splits an exactly degenerate shell is not caught by the curvature
alone.  The hydrogen halides with a frontier pair are the measured
case: the highest core orbital and the lowest active orbital are the
two components of a degenerate $\pi$ shell (core--active gap of order
$10^{-15}$~Ha), the energy is invariant to rotation inside that pair,
and the orbital Hessian acquires a zero direction whose sign is set by
rounding.  In the tested HCl/def2-SVP case one route reported zero
negative eigenvalues and the single-stage route two, at the same
space.  The gradient guard is then the only criterion that separates
such a window from a certified minimum.  Include both components of a
degenerate shell when a window is specified by hand; the automatic
fallback already does.

The curvature route is a separate matter.  At an exactly symmetric
stationary point --- reached by a manually specified N$_2$/6-31G
CAS(4,4) and by the automatically selected N quartet CAS(3,3) --- the
analytic orbital Hessian is not finite, while the energy and the
analytic gradient are.  The driver then takes the curvature from a
symmetric difference of the analytic gradient and records
\texttt{curvature\_method = 'gradient-difference'} in the result.

The fallback gives a \emph{starting} space, not a certified one.  It
says only what the occupations can say; whether the space holds is the
question the robust-level entanglement verification answers.
